\begin{SCn}
\begin{small}

\scnheader{NELL}
\scnidtf{Автоматизированная система непрерывного извлечения знаний из интернета, разработанная Карнеги-Меллонским университетом}
\scnidtf{Система извлечения знаний, разработанная в Питтсбурге}
\scniselement{инструмент для построения графов знаний}
\begin{scnindent}
    \scnidtf{система извлечения знаний для построения графов}
\end{scnindent}
\scniselement{семантическое программное обеспечение}
\scnrelfrom{разработчик}{Карнеги-Меллонский университет}
\scnrelfrom{год создания}{2010}
\scntext{назначение}{Автоматическое пополнение баз знаний через извлечение триплетов (сущность, отношение, сущность) из текстов веб-страниц}
\begin{scnrelfromset}{функциональные возможности}
    \scnfileitem{постоянное извлечение знаний}
    \begin{scnindent}
        \scnidtf{работа 24/7 и извлечение информации в виде сущностей и их взаимосвязей}
        \scntext{пример использования}{Анализ текстовых данных из интернета}
    \end{scnindent}
    \scnfileitem{создание онтологий}
    \begin{scnindent}
        \scnidtf{преобразование извлеченной информации в структурированный формат (например, RDF)}
        \scntext{пример использования}{Введение новой терминологии и правил формирования заголовков в библиографическом описании}
    \end{scnindent}
    \scnfileitem{масштабируемость и долгосрочная адаптация}
    \begin{scnindent}
        \scnidtf{улучшение точности извлечения информации со временем, благодаря адаптации к новым данным и расширения своих знаний без перезапуска}
        \scntext{пример использования}{Разработка методов компактного табличного представления функций, позволяющих эффективно обрабатывать данные в долгосрочной перспективе}
    \end{scnindent}
    \scnfileitem{семантическая обработка}
    \begin{scnindent}
        \scnidtf{фокусировка на понимании контекста и семантики текста, что способствует выявлению сложных взаимосвязей между сущностями}
        \scntext{пример использования}{Анализ контекста в задачах самоконтроля качества подготовки к занятиям, где важно понимание смысла вопросов и ответов}
    \end{scnindent}
\end{scnrelfromset}
\begin{scnrelfromset}{поддерживаемые языки}
    \scnitem{RDF}
    \scnitem{Turtle}
\end{scnrelfromset}
\begin{scnrelfromset}{преимущества}
    \scnfileitem{универсальность и модульность}
    \begin{scnindent}
        \scntext{уточнение}{Позволяет адаптировать функции под различные задачи, что делает её гибкой для применения в разных областях.}
    \end{scnindent}
    \scnfileitem{высокая эффективность}
    \scnfileitem{непрерывное обучение}
    \begin{scnindent}
        \scntext{уточнение}{Имитирует человеческий стиль обучения, позволяя системе постепенно накапливать знания и улучшать понимание языка через контекст.}
    \end{scnindent}
\end{scnrelfromset}
\begin{scnrelfromset}{недостатки}
    \scnfileitem{энергозависимость}
    \scnfileitem{Сложность реализации}
    \begin{scnindent}
        \scntext{уточнение}{Для эффективной работы требуется мощная инфраструктура и алгоритмы, способные обрабатывать большие объёмы данных в реальном времени.}
    \end{scnindent}
\end{scnrelfromset}
\begin{scnrelfromset}{библиографические источники}
    \scnfileitem{NELL: The Computer that Learns // Carnegie Mellon University. – Электронный ресурс. Дата обращения: 26.05.2025. https://www.cmu.edu/homepage/computing/2010/fall/nell-computer-that-learns.shtml}
    \scnfileitem{Never-Ending Language Learning (NELL) // Wikipedia. – Электронный ресурс. Дата обращения: 26.05.2025. https://en.wikipedia.org/wiki/NELL}
    \scnfileitem{RTW Resources // Carnegie Mellon University. – Электронный ресурс. Дата обращения: 26.05.2025. http://rtw.ml.cmu.edu/rtw/resources}
    \scnfileitem{NELL in Ontology Development // IBM Journal of Research and Development. – 2012. – Т. 56, № 3–4. – С. 12:1–12:10. – DOI: 10.1147/JRD.2012.2185930.}
    \scnfileitem{Never-Ending Learning / Carnegie Mellon University School of Computer Science. – Pittsburgh, PA, 2014. – 14 н. (Технический отчет). Доступ через: https://www.cs.cmu.edu/\~{}tom/pubs/NELL\_aaai15.pdf}
\end{scnrelfromset}

\scnheader{XLore}
\scnidtf{Англо-китайский билингвальный граф знаний на основе Wikipedia и Baidu Baike.}
\scniselement{крупномасштабный знаниевый граф}
\scniselement{семантическое программное обеспечение}
\scnrelfrom{разработчик}{Университет Цинхуа}
\scnrelfrom{год создания}{2013}
\scntext{назначение}{Инструмент предназначен для интеграции и структурирования межъязыковых знаний, предоставляя единое пространство для работы с англо-китайскими данными}
\begin{scnrelfromset}{архитектурные особенности}
    \scnfileitem{многоязычная интеграция}
    \scnfileitem{поддержка API}
    \scnfileitem{использование Python и JavaScript}
\end{scnrelfromset}
\begin{scnrelfromset}{функциональные возможности}
    \scnfileitem{поиск сущностей и концепций}
    \begin{scnindent}
        \scntext{уточнение}{Позволяет находить сущности, концепции и свойства на основе ключевых слов через API}
    \end{scnindent}
    \scnfileitem{межъязыковые связи}
    \scnfileitem{обработка естественного языка}
    \scnfileitem{автоматизация извлечения знаний}
    \begin{scnindent}
        \scntext{уточнение}{Платформа использует методы моделирования и извлечения знаний для автоматического построения графа из неструктурированных и полуструктурированных источников}
    \end{scnindent}
\end{scnrelfromset}
\begin{scnrelfromset}{преимущества}
    \scnfileitem{комплексная схема (онтология)}
    \begin{scnindent}
        \scntext{уточнение}{Используется тщательно спроектированная схема XLORE Ontology, которая эффективно организует и управляет сущностями из разнородных источников, обеспечивая структурированное представление знаний}
    \end{scnindent}
    \scnfileitem{многоязычная интеграция}
    \begin{scnindent}
        \scntext{уточнение}{Объединяет данные на китайском и английском языках, создавая сбалансированный кросс-лингвальный граф знаний с богатыми межъязыковыми связями, что полезно для мультиязыковых приложений}
    \end{scnindent}
    \scnfileitem{масштабируемость и объем данных}
    \begin{scnindent}
        \scntext{уточнение}{Включает в себя более 50 миллионов сущностей и 500 миллионов триплетов, а также предоставляет API с ежегодным доступом более 20 миллионов запросов}
    \end{scnindent}
    \scnfileitem{автоматированное построение знаний}
\end{scnrelfromset}
\begin{scnrelfromset}{недостатки}
    \scnfileitem{зависимость от ручного проектирования}
    \begin{scnindent}
        \scnexplanation{Некоторые методы построения графа требуют ручного создания шаблонов связей между сущностями, что ограничивает автоматизацию и увеличивает трудозатраты}
    \end{scnindent}
    \scnfileitem{ограниченная адаптация к большим данным}
    \begin{scnindent}
        \scnexplanation{Методы, использованные в XLore, могут сталкиваться с трудностями при обработке очень крупных наборов данных}
    \end{scnindent}
    \scnfileitem{структурные ограничения}
    \begin{scnindent}
        \scnexplanation{Граф может содержать недостаточно гибкие связи между сущностями, что затрудняет их использование в сложных аналитических задачах}
    \end{scnindent}
\end{scnrelfromset}
\begin{scnrelfromset}{библиографические источники}
    \scnfileitem{XLORE: Кросс-лингвистический граф знаний // XLORE. – Электронный ресурс. Дата обращения: 26.05.2025. https://xlore.cn/AboutUs}
    \scnfileitem{Тан, Цзя. Интеллектуальные системы будущего: Wu Dao – крупнейшая в мире модель ИИ / Цзя Тан // Tsinghua University. – Электронный ресурс – 2021. Дата обращения: 26.05.2025. https://www.tsinghua.edu.cn/en/info/1418/10451.htm}
    \scnfileitem{ACM Digital Library // DOI: 10.1145/3660521. – Электронный ресурс. Дата обращения: 26.05.2025. https://dl.acm.org/doi/full/10.1145/3660521}
    \scnfileitem{Научные данные // SciDB. – Электронный ресурс. Дата обращения: 26.05.2025. https://www.scidb.cn/en/detail?dataSetId=760439203967270912}
    \scnfileitem{XLORE: Крупномасштабный кросс-лингвистический граф знаний / Цзя Тан [и др.] // Proceedings of the International Semantic Web Conference (ISWC 2013). – CEUR-WS, 2013. – Электронный ресурс. Дата обращения: 26.05.2025. https://ceur-ws.org/Vol-1035/iswc2013\_demo\_31.pdf }
\end{scnrelfromset}

\bigskip

\scnheader{Сравнение редакторов онтологий Protégé и TopBraid Composer}
\begin{scneqtovector}
    \scnitem{NELL}
    \scnitem{XLoer}
\end{scneqtovector}
\begin{scnrelfromset}{сходства}
    \scnfileitem{Предназначены для извлечения информации из неструктурированных источников}
    \scnfileitem{Обеспечивают семантическую обработку данных для выявления сложных взаимосвязей между сущностями}
    \scnfileitem{Используются в научных и промышленных проектах для структурирования знаний}
    \scnfileitem{Позволяют масштабировать данные и адаптироваться к новым источникам информации}
\end{scnrelfromset}
\begin{scnrelfromset}{различия}
    \scnfileitem{NELL фокусируется на непрерывном извлечении триплетов (сущность, отношение, сущность) из интернета, в то время как XLore специализируется на межъязыковой интеграции англо-китайских данных}
    \scnfileitem{NELL требует мощной инфраструктуры для обработки больших данных в реальном времени, XLore сталкивается с ограничениями при обработке очень крупных наборов из-за зависимости от ручного проектирования}
    \scnfileitem{NELL ориентирован на автоматическое пополнение баз знаний без перезапуска, XLore требует частичной ручной корректировки для обеспечения гибкости связей.}
\end{scnrelfromset}
\begin{scnrelfromset}{рекомендации по выбору}
    \scnfileitem{NELL рекомендуется для проектов, требующих непрерывного извлечения знаний из интернета и долгосрочного обучения, например, в области анализа текстовых данных или создания баз знаний с контекстным пониманием.}
    \scnfileitem{XLore целесообразно использовать для задач, связанных с мультиязыковыми приложениями, такими как переводы, межъязыковые поисковые системы или интеграция данных из англоязычных и китайских источников.}
    \scnfileitem{NELL подходит для научных исследований, где важна автоматизация извлечения и обновления знаний, XLore — для корпоративных решений, требующих структурированного представления межъязыковых данных.}
    \scnfileitem{Если требуется высокая масштабируемость и обработка больших объёмов текста в реальном времени, предпочтение стоит отдать NELL. Для задач с акцентом на межъязыковую совместимость и API-интеграцию лучше выбрать XLore.}
\end{scnrelfromset}

\end{small}
\end{SCn}
